\documentclass[11pt]{article}
\usepackage[margin=1in]{geometry}
\usepackage{amsmath,amssymb,amsthm,mathtools}
\usepackage{booktabs}
\usepackage[colorlinks=true,linkcolor=blue,citecolor=blue,urlcolor=blue]{hyperref}

\newtheorem{theorem}{Theorem}
\newtheorem{corollary}{Corollary}
\newtheorem{proposition}{Proposition}
\newtheorem{lemma}{Lemma}
\theoremstyle{definition}
\newtheorem{definition}{Definition}
\theoremstyle{remark}
\newtheorem{remark}{Remark}

\newcommand{\ket}[1]{\lvert #1 \rangle}
\newcommand{\bra}[1]{\langle #1 \rvert}
\newcommand{\braket}[2]{\langle #1 \vert #2 \rangle}

\title{A Blindness Hierarchy for Quantum Benchmarks:\\
Conjugation Invariance and What Standard Metrics Cannot See}
\author{Cleiton Augusto Correa Bezerra\\
\small \texttt{augusto.cleiton@gmail.com}}
\date{July 2026 --- working notes, v1}

\begin{document}
\maketitle

\begin{abstract}
We prove that every quantum benchmark whose score is a function of
computational-basis measurement statistics of circuits prepared from real
initial states --- a class that includes Quantum Volume, linear
cross-entropy benchmarking (XEB), mirror circuit benchmarking, randomized
benchmarking, and the QED-C application suite --- is exactly blind to
\emph{consistent complex conjugation}: a simulator or device implementing
every gate $U$ as its complex conjugate $U^{*}$ achieves identical scores
to a correct one. The invariance extends to adaptive protocols with
mid-circuit measurement, and, strikingly, to full Quantum Phase
Estimation, which returns the \emph{correct} phase on a simulator in
which every gate is wrong. We then prove a relative-detectability
theorem for \emph{partial} conjugation bugs (the class found in
production in quantrs2 v0.2.0): such bugs are invisible on any circuit
family whose remaining gates are real --- which includes the standard
Quantum Volume $R_z$/$R_y$/CX model circuits and textbook QPE --- yet
are separated with probability gap $|\sin\theta|$ by a crafted four-gate
witness. All results are verified to machine precision by an open-source
Rust implementation. The practical conclusion: benchmark scores certify
sampling behaviour, not gate-convention correctness; closing the gap
requires amplitude-level, cross-implementation conformance testing.
\end{abstract}

\section{Setup and notation}

Fix $n$ qubits with computational basis
$\{\ket{x}\}_{x\in\{0,1\}^{n}}$. A \emph{circuit} is a finite word
$C = U_{L}\cdots U_{2}U_{1}$ of unitary gates. For a matrix $A$ written
in the computational basis, $A^{*}$ denotes entrywise complex
conjugation. We call a vector or matrix \emph{real} if it equals its
conjugate.

\begin{definition}[Conjugated circuit]
For a circuit $C = U_{L}\cdots U_{1}$, its conjugation is
$C^{*} \coloneqq U_{L}^{*}\cdots U_{1}^{*}$.
\end{definition}

Note that $(AB)^{*}=A^{*}B^{*}$, so conjugating every gate of a circuit
conjugates the composed unitary; the map $C \mapsto C^{*}$ is an
automorphism of the circuit monoid. Each $U^{*}$ is again unitary, so a
``conjugated simulator'' --- one whose gate library carries a consistent
$i \mapsto -i$ convention error --- simulates a perfectly valid quantum
evolution, namely the image of the intended one under the antiunitary
symmetry $K$ (complex conjugation in the computational basis):
$U^{*} = K U K^{-1}$.

\begin{definition}[Sampling benchmark]
A \emph{sampling benchmark} is any protocol whose score is a
(possibly randomized) function of the joint outcome statistics of
computational-basis measurements performed on circuits applied to the
initial state $\ket{0\cdots 0}$ (or any real initial state), possibly
adaptively and with mid-circuit measurements.
\end{definition}

This definition covers Quantum Volume heavy-output generation
\cite{cross2019}, linear XEB \cite{arute2019}, mirror circuit
benchmarking \cite{proctor2022}, randomized benchmarking in all its
standard variants, the QED-C application-oriented suite, and direct
algorithm tests (Grover, Bernstein--Vazirani, QPE) scored on output
histograms.

\section{The invariance theorem}

\begin{theorem}[Conjugation invariance]\label{thm:main}
Let $C$ be any circuit and $\ket{\psi_{0}}$ a real state. Then for every
basis outcome $x$,
\[
\bigl|\braket{x}{C^{*}\,\psi_{0}}\bigr|^{2}
= \bigl|\braket{x}{C\,\psi_{0}}\bigr|^{2}.
\]
\end{theorem}

\begin{proof}
Since $\ket{x}$ and $\ket{\psi_{0}}$ are real,
$\braket{x}{C^{*}\psi_{0}}
= \bra{x}^{*} C^{*} \ket{\psi_{0}}^{*}
= \bigl(\braket{x}{C\,\psi_{0}}\bigr)^{*}$,
and $|z^{*}|^{2}=|z|^{2}$.
\end{proof}

\begin{theorem}[Invariance for adaptive protocols]\label{thm:adaptive}
Any sampling benchmark (Definition~2), including protocols with
mid-circuit computational-basis measurements and classical feed-forward,
assigns identical outcome distributions --- and therefore identical
scores --- to the gate set $\{U_{g}\}$ and to $\{U_{g}^{*}\}$.
\end{theorem}

\begin{proof}
By induction over the protocol tree. Invariant: before each step, the
states reached under the two gate sets on any fixed transcript of past
outcomes are conjugates of one another, $\ket{\phi'} = \ket{\phi}^{*}$;
initially both equal the real $\ket{\psi_{0}}$. Unitary steps preserve
the invariant since $U^{*}\ket{\phi}^{*} = (U\ket{\phi})^{*}$.
A computational-basis measurement has real projectors
$P_{x} = \ket{x}\bra{x}$, so outcome probabilities agree,
$\bra{\phi^{*}}P_{x}\ket{\phi^{*}} = \bra{\phi}P_{x}\ket{\phi}$,
and post-measurement states remain conjugate:
$P_{x}\ket{\phi}^{*}/\|\cdot\| = \bigl(P_{x}\ket{\phi}/\|\cdot\|\bigr)^{*}$.
Classical feed-forward depends only on outcomes, which agree in
distribution. Hence the full joint outcome distribution is identical,
and so is any score computed from it.
\end{proof}

\begin{remark}[Physical interpretation]
Complex conjugation in a fixed basis is the canonical antiunitary
symmetry of quantum theory. Theorem~\ref{thm:adaptive} says that from
\emph{within} a single laboratory that prepares real states and measures
in the computational basis, the choice of $i$ versus $-i$ is a gauge:
no experiment of this form can resolve it. (Distinguishing quantum
theory from its real/conjugate images is known to require network
scenarios with independent sources \cite{renou2021}; our observation is
the operational, single-device special case --- and its consequence for
software testing, which appears to be unexamined.)
\end{remark}

\section{Corollaries: the standard benchmark suite is blind}

\begin{corollary}[Quantum Volume / HOG]\label{cor:qv}
The heavy-output-generation probability of a device is invariant under
consistent conjugation of its gate set. In particular a simulator
implementing $R_z(\theta)$ as $R_z(-\theta)$ passes Quantum Volume model
circuits (basis $R_z, R_y,$ CX) indistinguishably from a correct one,
since $R_y$ and CX are real and $R_z(-\theta) = R_z(\theta)^{*}$.
\end{corollary}

\begin{corollary}[Linear XEB]\label{cor:xeb}
The linear cross-entropy fidelity
$F_{\mathrm{XEB}} = 2^{n}\,\mathbb{E}_{x \sim p_{\mathrm{dev}}}
[\,p_{\mathrm{ideal}}(x)\,] - 1$
is invariant under conjugation of the device, of the ideal reference, or
both. \emph{The metric used to certify quantum computational advantage
does not distinguish $U$ from $U^{*}$.}
\end{corollary}

\begin{proof}
$p_{\mathrm{dev}}$ is invariant by Theorem~\ref{thm:main} applied to the
device; $p_{\mathrm{ideal}}$ is a probability of the ideal circuit from
$\ket{0}$, invariant for the same reason.
\end{proof}

\begin{corollary}[Mirror circuits]\label{cor:mirror}
Mirror benchmarking scores, being functions of the outcome statistics of
circuits $C^{\dagger} C$ (and Pauli-dressed variants) from
$\ket{0\cdots0}$, are conjugation-invariant. Directly:
$(C^{*})^{\dagger} C^{*} = (C^{\dagger} C)^{*}$ and
Theorem~\ref{thm:main} applies.
\end{corollary}

\begin{corollary}[Randomized benchmarking]\label{cor:rb}
Every RB variant whose data are survival frequencies of random gate
sequences measured in the computational basis is conjugation-invariant.
(No representation-theoretic argument is needed: RB is a sampling
benchmark in the sense of Definition~2, and Theorem~\ref{thm:adaptive}
applies verbatim. It is convenient that the Clifford group is closed
under conjugation, but the invariance does not depend on it.)
\end{corollary}

\begin{corollary}[Algorithm-level tests, including QPE]\label{cor:qpe}
Output histograms of Grover search, Bernstein--Vazirani, QFT-based
arithmetic, and full Quantum Phase Estimation are identical under
consistent conjugation. In particular QPE run entirely on a conjugated
simulator reports the \emph{correct} eigenphase: the conjugated
controlled-phase kickback ($\varphi \to -\varphi$) is exactly undone by
the conjugated inverse QFT, which rotates the other way.
\end{corollary}

\begin{remark}
Corollary~\ref{cor:qpe} refutes the natural conjecture that
phase-estimation algorithms act as conjugation detectors. They would
only if the QPE machinery (QFT, state preparation) ran on a
\emph{trusted correct} device while the gate under test carried the bug
--- a cross-implementation scenario, which is precisely the resolution
proposed in Section~\ref{sec:resolution}.
\end{remark}

\section{The hierarchy: partial bugs and relative detectability}

Real bugs are rarely library-wide. The quantrs2 v0.2.0 bug inverted the
rotation sign of $R_z$ alone. Call a bug \emph{$G$-partial} if exactly
the gates in a subset $G$ of the library are conjugated.

\begin{theorem}[Relative blindness of partial bugs]\label{thm:partial}
Let a device carry a $G$-partial conjugation bug, and let $\mathcal{F}$
be a circuit family in which every gate \emph{not} in $G$ is real. Then
on $\mathcal{F}$ the partial bug acts as a consistent global
conjugation, and every sampling benchmark built on $\mathcal{F}$ is
blind to it (Theorem~\ref{thm:adaptive}).
\end{theorem}

\begin{proof}
On circuits from $\mathcal{F}$, gates outside $G$ satisfy
$U = U^{*}$, so the buggy execution
$\cdots U^{*}_{g\in G} \cdots U_{h\notin G}\cdots$
coincides with $\cdots U^{*}_{g} \cdots U^{*}_{h}\cdots = C^{*}$.
\end{proof}

Theorem~\ref{thm:partial} explains why the quantrs2 bug survived every
standard benchmark: Quantum Volume model circuits use the basis
$\{R_z, R_y, \mathrm{CX}\}$ and textbook QPE uses
$\{H, X, \mathrm{SWAP}, \mathrm{controlled\text{-}phase}\}$; in both,
all non-phase-family gates are real, so the $R_z$-sign bug is
gauge-equivalent to a global conjugation on exactly the circuits these
benchmarks run.

Partial bugs, unlike global ones, \emph{are} detectable in principle ---
by breaking the realness condition deliberately:

\begin{proposition}[Four-gate witness]\label{prop:witness}
Consider a device with an $R_z$-sign bug but a correct $S$ gate. The
single-qubit circuit $H \cdot R_z(\theta)\cdot S \cdot H$ (gates listed
in application order) yields
\[
P(\,0\,) = \tfrac{1}{2}\bigl(1 - \sin\theta\bigr)
\quad\text{correct},
\qquad
P(\,0\,) = \tfrac{1}{2}\bigl(1 + \sin\theta\bigr)
\quad\text{buggy},
\]
a probability gap of $|\sin\theta|$ --- maximal ($\Delta p = 1$) at
$\theta = \pi/2$, and plainly visible in measurement counts.
\end{proposition}

\begin{proof}
$\ket{0} \xrightarrow{H} \tfrac{1}{\sqrt2}(\ket{0}+\ket{1})
\xrightarrow{R_z(\pm\theta)} \tfrac{1}{\sqrt2}
 (e^{\mp i\theta/2}\ket{0}+e^{\pm i\theta/2}\ket{1})
\xrightarrow{S} \tfrac{1}{\sqrt2}
 (e^{\mp i\theta/2}\ket{0}+e^{i(\pm\theta/2+\pi/2)}\ket{1})$,
so after the final $H$ the amplitude of $\ket{0}$ has squared modulus
$\tfrac14\,\lvert 1+e^{i(\pi/2\pm\theta)}\rvert^{2}
=\tfrac12\bigl(1+\cos(\pi/2\pm\theta)\bigr)
=\tfrac12\bigl(1\mp\sin\theta\bigr)$.
\end{proof}

The correct complex gate ($S$) plays the role of a fixed phase
reference; the theorem-level point is that \emph{some} correctly
implemented non-real operation must appear in the test circuit, or no
measurement statistics can see the bug. This is the design principle
behind phase-sensitive fuzzing: search the gate-mixing space that
benchmarks never enter.

The resulting hierarchy:

\begin{center}
\begin{tabular}{lll}
\toprule
Bug class & Detectable by sampling? & Detected by \\
\midrule
Global conjugation & Never (gauge symmetry) & amplitude-level cross-check \\
Partial conjugation & Only off the real-basis family & crafted witness / fuzzing \\
Non-conjugation errors & Generically yes & standard benchmarks \\
\bottomrule
\end{tabular}
\end{center}

\section{Empirical verification}

All claims were verified with the open-source CleitonForge toolkit
(Rust; \texttt{cforge-backends}), comparing a reference statevector
simulator against (i) a fully conjugated simulator and (ii) an
$R_z$/phase-family sign-bug simulator, on twenty $5$-qubit, depth-$5$
random $R_z/R_y/\mathrm{CX}$ circuits, $10^{5}$ shots:

\begin{center}
\begin{tabular}{lccc}
\toprule
Benchmark & correct & $U^{*}$ (all gates) & $R_z$-sign bug \\
\midrule
QV heavy-output prob. (seed 0) & $0.862400$ & $0.862400$ & $0.862400$ \\
Linear XEB, exact (seed 0) & $0.994064497$ & $0.994064497$ & $0.994064497$ \\
Mirror survival $P(\ket{0\cdots0})$ & $1.000000$ & $1.000000$ & $1.000000$ \\
QPE of $\varphi = 5/8$ (peak, prob.) & $5$ ($P{=}1$) & $5$ ($P{=}1$) & $5$ ($P{=}1$) \\
\midrule
Witness $H R_z(1) S H$: $P(0)$ & $0.079265$ & $0.079265$ & $0.920735$ \\
QGCS amplitude checks failed & $0/12$ & $4/12$ & $1/12$ \\
\bottomrule
\end{tabular}
\end{center}

Maximum benchmark-score deviation across all circuits and all three
devices: $0$ (bit-identical samples under a fixed seed, as the
distributions are equal to machine precision). The witness gap matches
Proposition~\ref{prop:witness} exactly:
$P(0)_{\mathrm{buggy}} - P(0)_{\mathrm{correct}} = \sin(1) = 0.841471$.

\section{Resolution: amplitude-level conformance}\label{sec:resolution}

The gauge is fixed by any procedure with access to \emph{relative
amplitudes} rather than probabilities alone --- operationally, by
comparing two implementations against an analytically derived amplitude
convention. The QGCS conformance suite does this with twelve
$O(1)$-size witness circuits checking amplitude \emph{ratios} (e.g.\
$R_z(\pi/2)\ket{+}$ must produce ratio $+i$ between its components
under the OpenQASM~3 convention); it separates all three devices above
instantly. The same principle generalizes to differential fuzzing with
an amplitude-equality oracle, which searches automatically for the
gate mixtures of Proposition~\ref{prop:witness}; this is developed in a
companion software-engineering paper.

\begin{thebibliography}{9}
\bibitem{cross2019} A.~Cross, L.~Bishop, S.~Sheldon, P.~Nation,
J.~Gambetta, \emph{Validating quantum computers using randomized model
circuits}, Phys.\ Rev.\ A \textbf{100}, 032328 (2019).
\bibitem{arute2019} F.~Arute \emph{et al.}, \emph{Quantum supremacy
using a programmable superconducting processor}, Nature \textbf{574},
505 (2019).
\bibitem{proctor2022} T.~Proctor \emph{et al.}, \emph{Measuring the
capabilities of quantum computers}, Nat.\ Phys.\ \textbf{18}, 75 (2022).
\bibitem{renou2021} M.-O.~Renou \emph{et al.}, \emph{Quantum theory
based on real numbers can be experimentally falsified}, Nature
\textbf{600}, 625 (2021).
\end{thebibliography}

\end{document}
